\documentclass[10pt,a4paper]{article}
\usepackage[T1]{fontenc}
\usepackage[utf8]{inputenc}
\usepackage[portuguese]{babel}
\usepackage[margin=1.6cm]{geometry}
\usepackage{enumitem}
\usepackage{titlesec}
\usepackage{xcolor}
\usepackage[hidelinks]{hyperref}
\usepackage{booktabs}

\definecolor{cinza}{gray}{0.35}
\titleformat{\section}{\normalsize\bfseries}{}{0pt}{}[\vspace{-5pt}\rule{\textwidth}{0.4pt}]
\titlespacing{\section}{0pt}{6pt}{3pt}
\newcommand{\lin}[2]{\noindent\textbf{#1}\hfill{\color{cinza}\small #2}\par}
\setlist[itemize]{leftmargin=1.1em,itemsep=1pt,topsep=1pt,parsep=0pt}
\pagestyle{empty}

\begin{document}

% Proprietário · Copyright (c) 2026 Aarão Melo Lopes

\begin{center}
{\LARGE\bfseries Aarão Melo Lopes}\\[1pt]
{\color{cinza}\small Engenheiro Eletricista \,|\, Sistemas \,|\, Matemática Aplicada}\\[1pt]
{\footnotesize
\href{mailto:aarao.melo.lopes@gmail.com}{aarao.melo.lopes@gmail.com}
\,·\,
\href{https://goldenkingdom.patriatechnology.com}{goldenkingdom.patriatechnology.com}
\,·\,
\href{https://goldenkingdom.patriatechnology.com/repo.git}{goldenkingdom\ldots/repo.git}}
\end{center}

\section*{Perfil}

Engenheiro eletricista (UFRR) e mestrado em Ciência da Computação (IME-USP, créditos).
Construo software de sistemas em C e a documentação que o verifica.
Procuro colaboração em laboratório, indústria ou engenharia de software.

\section*{Formação}

\lin{Mestrado em Ciência da Computação --- IME-USP}{2020--2022}
Bolsista CAPES. \textbf{72 créditos} (exigência: 44); conceito máximo em 8/9 disciplinas.
Trabalho teórico e sistemas no lugar da dissertação formal --- disponível sob consulta.

\smallskip
\lin{Bacharelado em Engenharia Elétrica --- UFRR}{2011--2018}
Média \textbf{8{,}66}. TCC nota 10 (controlo de motor / inversores), recomendação de publicação.
Monitoria: \textbf{473\,h} (Cálculo e Álgebra Linear).

\section*{Competências}

{\small C · Makefile · Shell · POSIX · Git · assembly ·
armazenamento persistente em disco · máquinas virtuais ·
testes automatizados · \textbf{testes por reversão e mutação} · documentação técnica}

\section*{Projetos em destaque}

\begin{itemize}
\item \textbf{Bateria de verificação} --- 328 testes em C a zero; CI bloqueia publicação se falhar.
Mede-se por \emph{reversão} --- resíduo zero --- e não por comparação contra valor esperado; e pelas
\textbf{duas metades}: o resíduo onde tem de ser zero e o resíduo onde não pode ser.
\item \textbf{Máquina de uma instrução} --- VM com ISA mínima de uma única instrução,
sem estado auxiliar (sem \texttt{.bss}).
\item \textbf{Orquestração de agentes} --- \texttt{broca-so}: runtime fixo (imutável em
execução) e passagem de estado entre agentes.
\item \textbf{Motor de xadrez + compilador} --- sistemas próprios, repositório público.
\end{itemize}

\section*{Software desenvolvido}

Compilador multi-formato, VM de pilha, núcleo mínimo, banco persistente em disco sem
\emph{malloc}, assistente local sem modelo externo. Stack: C, POSIX, Git.

\smallskip
\noindent \textbf{Dual Sort} --- ordenação que corre \emph{no disco}, sem memória dinâmica e sem
uma comparação: a ordem lê-se em vez de se produzir. Guardando o caminho, a operação é reversível e
\textbf{não apaga um bit} --- $0$ contra $25{,}6$ milhões na forma destrutiva, medido.

\smallskip
\noindent \textbf{Ordenar e cortar são duais}, e daí o \emph{corte máximo} (\emph{max-cut}) numa
árvore sai por leitura em vez de busca: a paridade da profundidade corta todas as arestas, que é o
máximo possível. O corte constrói-se rodando a sequência natural meia volta --- sem comparações e
sem escolher centro.

\section*{Pesquisa (documentação pública)}

Dois artigos auto-contidos --- \emph{O corpo estelar} (19\,pp.), a especificação do sistema, e
\emph{Dual Sort} (7\,pp.), a ordenação sem comparações. E três volumes abertos
(Teoria 61\,pp.\ · Catálogo 457\,pp.\ · narrativa 365\,pp.):
álgebra com prova escrita e teste executável por afirmação.
\\[1pt]
{\footnotesize\color{cinza}Ainda sem adopção externa nem citações independentes.}

\section*{O que procuro}

Laboratório ou indústria (materialização de interface / compósito modelado);
engenharia de software (agentes em produção); equipas de jogo/cinema (mundo já escrito).

\smallskip
\begin{center}\footnotesize
\href{https://goldenkingdom.patriatechnology.com/docs/}{documentação}
\,·\,
\href{https://goldenkingdom.patriatechnology.com/repo.git}{https://goldenkingdom.patriatechnology.com/repo.git}
\,·\,
Proprietário — contrato e pagamento
\end{center}

\end{document}
